\chapter{Confession of a Man of Wealth}

\lettrine{T}{he} Theatin entered deliberately, without being too much  astonished at the noise and agitation which anxiety for the cardinal's  health had raised in his household. <Come in, my reverend father,>  said Mazarin, after a last look at the ruelle, <come in and console  me.>

<That is my duty, my lord,> replied the Theatin.

<Begin by sitting down, and making yourself comfortable, for I am going  to begin with a general confession; you will afterwards give me a good  absolution, and I shall believe myself more tranquil.>

<My lord,> said the father, <you are not so ill as to make a  general confession urgent—and it will be very fatiguing—take  care.>

<You suspect, then, that it may be long, father?>

<How can I think it otherwise, when a man has lived so completely as your  eminence has done?>

<Ah! that is true!—yes—the recital may be long.>

<The mercy of God is great,> snuffled the Theatin.

<Stop,> said Mazarin; <there I begin to terrify myself with  having allowed so many things to pass which the Lord might reprove.>

<Is that not always so?> said the Theatin naively, removing further  from the lamp his thin pointed face, like that of a mole. <Sinners are so  forgetful beforehand, and scrupulous when it is too late.>

<Sinners?> replied Mazarin. <Do you use that word ironically,  and to reproach me with all the genealogies I have allowed to be made on my  account—I—the son of a fisherman, in fact?>\footnote{This is quite untranslatable—it being a play upon the words \textit{pècheur}  (with a grave over the first e), a sinner, and \textit{pêcheur} (with an accent circumflex over the first e), a fisherman.—TRANS.}

<Hum!> said the Theatin.

<That is a first sin, father; for I have allowed myself made to descend  from two old Roman consuls, S. Geganius Macerinus 1st, Macerinus 2d, and  Proculus Macerinus 3d, of whom the Chronicle of Haolander speaks. From  Macerinus to Mazarin the proximity was tempting. Macerinus, a diminutive, means  leanish, poorish, out of case. Oh! reverend father! Mazarini may now be carried  to the augmentative Maigre, thin as Lazarus. Look!>—and he showed  his fleshless arms.

<In your having been born of a family of fishermen I see nothing  injurious to you; for—St Peter was a fisherman; and if you are a prince  of the church, my lord, he was the supreme head of it. Pass on, if you  please.>

<So much the more for my having threatened with the Bastille a certain  Bounet, a priest of Avignon, who wanted to publish a genealogy of the Casa  Mazarini much too marvellous.>

<To be probable?> replied the Theatin.

<Oh! if I had acted up to his idea, father, that would have been the vice  of pride—another sin.>

<It was an excess of wit, and a person is not to be reproached with such  sorts of abuses. Pass on, pass on!>

<I was all pride. Look you, father, I will endeavour to divide that into  capital sins.>

<I like divisions, when well made.>

<I am glad of that. You must know that in 1630—alas! that is  thirty-one years ago\longdash>

<You were then twenty-nine years old, monseigneur.>

<A hot-headed age. I was then something of a soldier, and I threw myself  at Casal into the arquebusades, to show that I rode on horseback as well as an  officer. It is true, I restored peace between the French and the Spaniards.  That redeems my sin a little.>

<I see no sin in being able to ride well on horseback,> said the  Theatin; <that is in perfect good taste, and does honour to our gown. As a  Christian, I approve of your having prevented the effusion of blood; as a monk,  I am proud of the bravery a monk has exhibited.>

Mazarin bowed his head humbly. <Yes,> said he, <but the  consequences?>

<What consequences?>

<Eh! that damned sin of pride has roots without end. From the time that I  threw myself in that manner between two armies, that I had smelt powder and  faced lines of soldiers, I have held generals a little in contempt.>

<Ah!> said the father.

<There is the evil; so that I have not found one endurable since that  time.>

<The fact is,> said the Theatin, <that the generals we have  had have not been remarkable.>

<Oh!> cried Mazarin, <there was Monsieur le Prince. I have  tormented him thoroughly!>

<He is not much to be pitied: he has acquired sufficient glory, and  sufficient wealth.>

<That may be, for Monsieur le Prince; but M. Beaufort, for  example—whom I held suffering so long in the dungeon of Vincennes?>

<Ah! but he was a rebel, and the safety of the state required that you  should make a sacrifice. Pass on!>

<I believe I have exhausted pride. There is another sin which I am afraid  to qualify.>

<I can qualify it myself. Tell it.>

<A great sin, reverend father!>

<We shall judge, monseigneur.>

<You cannot fail to have heard of certain relations which I have  had—with her majesty the queen-mother;—the malevolent\longdash>

<The malevolent, my lord, are fools. Was it not necessary for the good of  the state and the interests of the young king, that you should live in good  intelligence with the queen? Pass on, pass on!>

<I assure you,> said Mazarin, <you remove a terrible weight  from my breast.>

<These are all trifles!—look for something serious.>

<I have had much ambition, father.>

<That is the march of great minds and things, my lord.>

<Even the longing for the tiara?>

<To be pope is to be the first of Christians. Why should you not desire  that?>

<It has been printed that, to gain that object, I had sold Cambria to the  Spaniards.>

<You have, perhaps, yourself written pamphlets without severely  persecuting pamphleteers.>

<Then, reverend father, I have truly a clean breast. I feel nothing  remaining but slight peccadilloes.>

<What are they?>

<Play.>

<That is rather worldly: but you were obliged by the duties of greatness  to keep a good house.>

<I like to win.>

<No player plays to lose.>

<I cheated a little.>

<You took your advantage. Pass on.>

<Well! reverend father, I feel nothing else upon my conscience. Give me  absolution, and my soul will be able, when God shall please to call it, to  mount without obstacle to the throne\longdash>

The Theatin moved neither his arms nor his lips. <What are you waiting  for, father?> said Mazarin.

<I am waiting for the end.>

<The end of what?>

<Of the confession, monsieur.>

<But I have ended.>

<Oh, no; your eminence is mistaken.>

<Not that I know of.>

<Search diligently.>

<I have searched as well as possible.>

<Then I shall assist your memory.>

<Do.>

The Theatin coughed several times. <You have said nothing of avarice,  another capital sin, nor of those millions,> said he.

<What millions, father?>

<Why, those you possess, my lord.>

<Father, that money is mine, why should I speak to you about that?>

<Because, you see, our opinions differ. You say that money is yours,  whilst I—I believe it is rather the property of others.>

Mazarin lifted his cold hand to his brow, which was beaded with perspiration.  <How so?> stammered he.

<This way. Your excellency had gained much wealth—in the service of  the king.>

<Hum! much—that is, not too much.>

<Whatever it may be, whence came that wealth?>

<From the state.>

<The state; that is the king.>

<But what do you conclude from that, father?> said Mazarin, who  began to tremble.

<I cannot conclude without seeing a list of the riches you possess. Let  us reckon a little, if you please. You have the bishopric of Metz?>

<Yes.>

<The Abbeys of St Clement, St Arnould, and St Vincent, all at  Metz?>

<Yes.>

<You have the Abbey of St Denis, in France, magnificent property?>

<Yes, father.>

<You have the Abbey of Cluny, which is rich?>

<I have.>

<That of St Medard at Soissons, with a revenue of one hundred thousand  livres?>

<I cannot deny it.>

<That of St Victor, at Marseilles,—one of the best in the  south?>

<Yes father.>

<A good million a year. With the emoluments of the cardinalship and the  ministry, I say too little when I say two millions a year.>

<Eh!>

<In ten years that is twenty millions—and twenty millions put out  at fifty per cent. give, by progression, twenty-three millions in ten  years.>

<How well you reckon for a Theatin!>

<Since your eminence placed our order in the convent we occupy, near St  Germain des Pres, in 1644, I have kept the accounts of the society.>

<And mine likewise, apparently, father.>

<One ought to know a little of everything, my lord.>

<Very well. Conclude, at present.>

<I conclude that your baggage is too heavy to allow you to pass through  the gates of Paradise.>

<Shall I be damned?>

<If you do not make restitution, yes.>

Mazarin uttered a piteous cry. <Restitution!—but to whom, good  God?>

<To the owner of that money,—to the king.>

<But the king did not give it all to me.>

<One moment,—does not the king sign the ordonances?>

Mazarin passed from sighs to groans. <Absolution! absolution!>  cried he.

<Impossible, my lord. Restitution! restitution!> replied the  Theatin.

<But you absolve me from all other sins, why not from that?>

<Because,> replied the father, <to absolve you for that  motive would be a sin for which the king would never absolve me, my  lord.>

Thereupon the confessor quitted his penitent with an air full of compunction.  He then went out in the same manner he had entered.

<Oh, good God!> groaned the cardinal. <Come here, Colbert, I  am very, very ill indeed, my friend.>  